\section*{Contoh}
\addcontentsline{toc}{section}{Contoh}

\begin{contoh}[Menelusuri DFA]
Gunakan DFA pada \figref{fig:ch06-dfa-genap} untuk memeriksa string
\(abba\). State awal \(q_0\). Membaca \(a\): pindah ke \(q_1\).
Membaca \(b\): tetap di \(q_1\). Membaca \(b\): tetap di \(q_1\).
Membaca \(a\): pindah ke \(q_0\). State akhir \(q_0\) adalah state
penerima, sehingga \(abba\) diterima. Jumlah \(a\) pada string adalah
dua (genap), sesuai definisi bahasa.
\end{contoh}

\begin{contoh}[Epsilon closure bertingkat]
Misalkan NFA memiliki transisi
\(q_0 \xrightarrow{\varepsilon} q_1\),
\(q_1 \xrightarrow{\varepsilon} q_2\),
\(q_2 \xrightarrow{a} q_3\),
\(q_3 \xrightarrow{\varepsilon} q_4\).
Maka
\[
\varepsilon\text{-closure}(\{q_0\}) = \{q_0,q_1,q_2\}.
\]
Setelah membaca \(a\), himpunan state menjadi
\(\{q_3\}\), lalu epsilon closure-nya menjadi
\(\{q_3,q_4\}\). Jika \(q_4\) adalah state penerima, string \(a\)
diterima.
\end{contoh}

\begin{contoh}[Subset construction]
Untuk NFA dengan transisi
\(q_0 \xrightarrow{a} \{q_0,q_1\}\),
\(q_0 \xrightarrow{b} \{q_0\}\),
\(q_1 \xrightarrow{b} \{q_2\}\), dan \(q_2\) penerima:
\begin{center}
\small
\begin{tabular}{@{}c l c c c@{}}
\toprule
\textbf{State DFA} & \textbf{Himpunan NFA} &
\textbf{\(a\)} & \textbf{\(b\)} & \textbf{Penerima} \\
\midrule
\(S_0\) & \(\{q_0\}\) & \(S_1\) & \(S_0\) & tidak \\
\(S_1\) & \(\{q_0,q_1\}\) & \(S_1\) & \(S_2\) & tidak \\
\(S_2\) & \(\{q_0,q_2\}\) & \(S_1\) & \(S_0\) & ya \\
\bottomrule
\end{tabular}
\end{center}
State penerima adalah \(S_2\), karena memuat \(q_2\). DFA hasilnya
memiliki tiga state.
\end{contoh}

\begin{contoh}[Minimalisasi]
Diberikan DFA dengan state \(\{0,1,2,3\}\), state penerima
\(\{3\}\), dan transisi:
\[
0 \xrightarrow{a} 1,\quad
0 \xrightarrow{b} 2,\quad
1 \xrightarrow{a} 1,\quad
1 \xrightarrow{b} 3,\quad
2 \xrightarrow{a} 1,\quad
2 \xrightarrow{b} 2,\quad
3 \xrightarrow{a} 1,\quad
3 \xrightarrow{b} 2.
\]
Partisi awal \(\{0,1,2\}\) dan \(\{3\}\). Tabel berikut mencatat
setiap iterasi.

\begin{center}
\small
\begin{tabular}{@{}c l l@{}}
\toprule
\textbf{Iterasi} & \textbf{Partisi} & \textbf{Alasan} \\
\midrule
0 & \(\{0,1,2\}, \{3\}\) & penerima versus bukan penerima \\
1 & \(\{0,1\}, \{2\}, \{3\}\) &
\(1 \xrightarrow{b} 3\), sedangkan \(2 \xrightarrow{b} 2\) \\
2 & \(\{0\}, \{1\}, \{2\}, \{3\}\) &
\(0 \xrightarrow{b} 2\), sedangkan \(1 \xrightarrow{b} 3\) \\
\bottomrule
\end{tabular}
\end{center}

Hasil akhirnya empat grup; DFA sudah minimal.
\end{contoh}

\begin{contoh}[Token recognition]
Misalkan DFA gabungan memiliki state penerima untuk identifier dan
angka. Pada masukan \code{abc123}, DFA membaca seluruh string dan
berhenti di state penerima identifier. Pada masukan \code{123abc},
DFA membaca \code{123} sebagai angka, lalu berhenti karena transisi
berikutnya tidak ada. Lexer mengeluarkan token angka \code{123} dan
memproses \code{abc} pada pemanggilan berikutnya. Maximal munch
memastikan \code{123} tidak dipecah menjadi tiga token digit.
\end{contoh}
